\documentclass[11pt,a4paper]{article}

\usepackage{amsmath,amssymb,amsthm}
\usepackage{mathtools}
\usepackage{hyperref}
\usepackage[margin=2.5cm]{geometry}
\usepackage{enumitem}
\usepackage{graphicx}
\usepackage[normalem]{ulem}

\newtheorem{theorem}{Theorem}[section]
\newtheorem{lemma}[theorem]{Lemma}
\newtheorem{proposition}[theorem]{Proposition}
\newtheorem{corollary}[theorem]{Corollary}
\newtheorem{remark}[theorem]{Remark}
\theoremstyle{definition}
\newtheorem{definition}[theorem]{Definition}
\newtheorem{fact}[theorem]{Fact}

\DeclareMathOperator{\re}{Re}
\DeclareMathOperator{\im}{Im}
\newcommand{\D}{\mathbb{D}}
\newcommand{\R}{\mathbb{R}}
\newcommand{\C}{\mathbb{C}}
\newcommand{\N}{\mathbb{N}}
\newcommand{\T}{\mathbb{T}}

\title{Pair-Bound Lyapunov Stability for a Scalar Dissipative\\
Kuramoto-Type Mean-Field Equation, Formalized in Lean~4}
\author{}
\date{}

\emergencystretch=1.5em

\begin{document}
\maketitle

\begin{abstract}
We establish global asymptotic stability of the partially locked state for a scalar dissipative mean-field equation that shares the self-consistency structure of the Kuramoto--Ott--Antonsen system. The central device is a \emph{pair bound}: a pointwise algebraic inequality, lifted to the continuum via Fubini, that renders the $L^2$ Lyapunov function non-increasing without recourse to spectral theory or Fourier analysis. Combining this with a dominated-convergence bootstrap that prevents the order parameter from vanishing, we obtain convergence $r(t)\to r^*$ from \emph{all} initial data---no basin-of-attraction assumption is needed.

The entire argument is formalized in Lean~4 (Mathlib v4.30.0-rc1) with no \texttt{sorry} and no project-specific axioms. A conditional extension to the standard complex Ott--Antonsen equation and the full Kuramoto PDE is also provided, relying on two cited external results (nonlinear Landau damping and OA manifold attractivity). To the best of our knowledge, this constitutes the first machine-checked proof of a stability theorem for a Kuramoto-type model.
\end{abstract}

\paragraph{Main contributions.}
\begin{enumerate}[label=(\arabic*)]
\item A novel \emph{pair bound} (Proposition~\ref{prop:pair}) giving $\dot V \leq 0$ for the scalar dissipative model~\eqref{eq:oa}, with no antecedent in the Kuramoto literature.
\item Global stability without a basin condition: a DCT bootstrap (Proposition~\ref{prop:rpos}) shows $r(t) \geq r_{\min} > 0$ for all $K > K_c$, eliminating the assumption $V(0) < (r^*)^2$.
\item A complete machine-checked formalization covering Gaussian, Student-$t$ ($\nu>1$), compactly supported, and mixture distributions satisfying $\int\gamma\,\mathrm{d}\mu<\infty$.
\item A conditional bridge to the standard complex OA and full Kuramoto PDE via rotation cancellation and two external axioms.
\end{enumerate}

\tableofcontents

% ============================================================
\section{Setup and statement of results}\label{sec:setup}
% ============================================================

\subsection{The standard complex Ott--Antonsen equation}\label{sec:complex-oa}

On the Ott--Antonsen manifold, the Kuramoto--Sakaguchi PDE reduces to
\begin{equation}\label{eq:complex-oa}
\dot{z}(\omega,t) = -i\omega\,z + \frac{K}{2}\bigl(\bar\eta - \eta\,z^2\bigr), \qquad \eta(t) = \int_\R \overline{z(\omega,t)}\,g(\omega)\,\mathrm{d}\omega,
\end{equation}
where $z(\omega,t)\in\D$ and $\eta(t)\in\C$ is the complex order parameter~\cite{ott-antonsen2008}. When $g$ is symmetric and the initial data satisfy $z(\omega,0) = \overline{z(-\omega,0)}$, this symmetry persists and $\eta(t) = r(t)\in\R$. For the Lorentzian $g(\omega) = \gamma/(\pi(\omega^2+\gamma^2))$, residue calculus collapses~\eqref{eq:complex-oa} to the scalar ODE $\dot r = (K/2-\gamma)r - (K/2)r^3$.

\subsection{A scalar dissipative model}\label{sec:scalar-model}

We study the scalar real mean-field equation
\begin{equation}\label{eq:oa}
\dot{\alpha}(\omega,t) = -\gamma(\omega)\,\alpha + \frac{K}{2}\,r(t)\bigl(1 - \alpha^2\bigr), \qquad r(t) = \int_\R \alpha(\omega,t)\,g(\omega)\,\mathrm{d}\omega,
\end{equation}
with $\alpha(\omega,t)\in(0,1)$ and dissipation rate $\gamma(\omega)>0$. We write $\mathrm{d}\mu = g\,\mathrm{d}\omega$ for the probability measure induced by the density~$g$.

This is a \emph{distinct} dynamical system from~\eqref{eq:complex-oa}: the dissipative term $-\gamma\alpha$ replaces the rotational term $-i\omega z$. What the two systems share is the self-consistency equation defining the partially locked state.

\begin{remark}[Scope]\label{rem:scope}
Equation~\eqref{eq:oa} is \textbf{not} obtained from~\eqref{eq:complex-oa} by restricting to real~$z$; the subspace $\im(z)=0$ is not invariant under~\eqref{eq:complex-oa}. The Lyapunov argument in Sections~\ref{sec:proof}--\ref{sec:cauchy} applies to~\eqref{eq:oa} directly. The conditional extension to~\eqref{eq:complex-oa} appears in Section~\ref{sec:complex-extension}.
\end{remark}

\begin{center}\small
\begin{tabular}{@{}lllll@{}}
\textbf{Model} & \textbf{Eq.} & \textbf{Uncond.?} & \textbf{Lean status} & \textbf{Relation to Kuramoto} \\
\hline
Scalar dissipative & \eqref{eq:oa} & Yes & complete & Shares PLS self-consistency \\
Lorentzian ODE & Bernoulli & Yes & complete & Classical OA for Lorentzian $g$ \\
$n$-pole & Finite & Yes & complete & Discrete analogue \\
Complex OA & \eqref{eq:complex-oa} & No (1 axiom) & complete & Standard OA reduction \\
Full Kuramoto PDE & Kinetic & No (2 axioms) & complete & Original model \\
\end{tabular}
\end{center}

\subsection{The partially locked state (PLS)}

We assume the existence of a unique positive equilibrium $(\alpha^*,r^*)$ with $r^*\in(0,1)$ satisfying
\begin{equation}\label{eq:sc}
r^* = \int_\R \alpha^*(\omega)\,g(\omega)\,\mathrm{d}\omega
\end{equation}
and the pointwise balance
\begin{equation}\label{eq:equil}
\gamma(\omega)\,\alpha^*(\omega) = \frac{K}{2}\,r^*\bigl(1 - \alpha^*(\omega)^2\bigr) \qquad \text{for $g$-a.e.\ $\omega$.}
\end{equation}
For the Lorentzian this holds whenever $K>K_c := 2\gamma$.

\subsection{Main theorems}

\begin{theorem}[Basin stability]\label{thm:main}
Let $\gamma>0$ with $\int\gamma\,\mathrm{d}\mu<\infty$. Let $(\alpha^*,r^*)$ satisfy~\eqref{eq:equil}--\eqref{eq:sc}, and let $(\alpha,r)$ solve~\eqref{eq:oa} with:
\begin{enumerate}[label=(\roman*)]
\item $r(t) = \int\alpha(\omega,t)\,\mathrm{d}\mu$ and $\alpha(\omega,t)\in(0,1)$ for all $\omega,t\geq 0$;
\item $V(0) := \int(\alpha(\omega,0)-\alpha^*(\omega))^2\,\mathrm{d}\mu < (r^*)^2$;
\item for each $M>0$, there exists $\delta_0>0$ with $\alpha(\omega,0)\geq\delta_0$ on $\{\gamma\leq M\}$.
\end{enumerate}
Then $r(t)\to r^*$ as $t\to+\infty$.
\end{theorem}

\begin{remark}
No persistence hypothesis appears among the assumptions. The proof \emph{derives} order-parameter persistence ($r(t)\geq r_{\min}>0$) from $V$~antitonicity and Cauchy--Schwarz, then body persistence from the ODE barrier. Conditions~(ii)--(iii) hold for any initial datum in a neighborhood of the PLS.
\end{remark}

\begin{theorem}[Global stability]\label{thm:global}
Assume additionally that $\int(1/\gamma)\,\mathrm{d}\mu<\infty$ and $K > K_c := 2/\!\int\!(1/\gamma)\,\mathrm{d}\mu$. Then the conclusion of Theorem~\ref{thm:main} holds \textbf{without} condition~(ii): $r(t)\to r^*$ for all initial data satisfying~(i) and~(iii).
\end{theorem}

\begin{remark}
The threshold $K>K_c$ does not alone guarantee existence of $(\alpha^*,r^*)$; that is an additional hypothesis, verified for specific families (Gaussian, Student-$t$) via the implicit function theorem on~\eqref{eq:sc}.
\end{remark}

% ============================================================
\section{Proof of basin stability}\label{sec:proof}
% ============================================================

\subsection{Step 1: Leibniz rule}\label{sec:leibniz}

Define the $L^2$ Lyapunov function
\begin{equation}\label{eq:V}
V(t) := \int_\R (\alpha(\omega,t) - \alpha^*(\omega))^2\,\mathrm{d}\mu.
\end{equation}

\begin{proposition}[Leibniz]\label{prop:leibniz}
Under the hypotheses of Theorem~\ref{thm:main}, $V$ is continuous and
\[
V'(t) = \int_\R 2\,(\alpha - \alpha^*)\,\dot\alpha\,\mathrm{d}\mu \quad\text{for all } t>0.
\]
\end{proposition}

\begin{proof}
The pointwise derivative $\partial_t(\alpha-\alpha^*)^2 = 2(\alpha-\alpha^*)\dot\alpha$ satisfies $|2(\alpha-\alpha^*)\dot\alpha| \leq 2(\gamma(\omega)+K)$, since $|\alpha-\alpha^*|\leq 1$ and $|\dot\alpha|\leq\gamma+K$ on $(0,1)$. The right-hand side is $\mu$-integrable by hypothesis, so differentiation under the integral sign applies.
\end{proof}

\subsection{Step 2: The pair bound}\label{sec:pair}

\begin{proposition}[Pair bound]\label{prop:pair}
For any $\alpha_j,\alpha_k,\alpha^*_j,\alpha^*_k\in(0,1)$:
\begin{equation}\label{eq:pair}
\alpha^*_j(\alpha_k-\alpha^*_k)^2\!\Bigl(\alpha_k+\tfrac{1}{\alpha^*_k}\Bigr) + \alpha^*_k(\alpha_j-\alpha^*_j)^2\!\Bigl(\alpha_j+\tfrac{1}{\alpha^*_j}\Bigr) \geq (\alpha_j-\alpha^*_j)(\alpha_k-\alpha^*_k)(2-\alpha_j^2-\alpha_k^2).
\end{equation}
\end{proposition}

\begin{proof}
Write $p_j=\alpha_j-\alpha^*_j$, $p_k=\alpha_k-\alpha^*_k$. Multiplying both sides by $\alpha^*_j\alpha^*_k>0$ and expanding:
\begin{equation}\label{eq:decomp}
\alpha^*_j\alpha^*_k(\text{LHS}-\text{RHS}) = \underbrace{(\alpha^*_jp_k-\alpha^*_kp_j)^2}_{T_1} + \underbrace{(\alpha^*_j)^2\alpha_k\alpha^*_kp_k^2}_{T_2} + \underbrace{(\alpha^*_k)^2\alpha_j\alpha^*_jp_j^2}_{T_3} + \underbrace{\alpha^*_j\alpha^*_k(\alpha_j^2+\alpha_k^2)p_jp_k}_{T_4}.
\end{equation}

When $p_jp_k\leq 0$: each LHS summand is non-negative (factors involving $\alpha+1/\alpha^*$ are positive) while the RHS is non-positive (since $2-\alpha_j^2-\alpha_k^2>0$).

When $p_jp_k>0$: the first three terms of~\eqref{eq:decomp} are manifestly non-negative. The fourth is positive since $\alpha_j^2+\alpha_k^2>0$ and $p_jp_k>0$, so $\alpha_j=\alpha^*_j+p_j>0$ forces every factor positive.
\end{proof}

\begin{proposition}[Continuum pair bound]\label{prop:cpair}
The double integral
\[
\iint \Bigl[\alpha^*_1(\alpha_2-\alpha^*_2)^2\!\bigl(\alpha_2+\tfrac{1}{\alpha^*_2}\bigr) + \alpha^*_2(\alpha_1-\alpha^*_1)^2\!\bigl(\alpha_1+\tfrac{1}{\alpha^*_1}\bigr) - (\alpha_1-\alpha^*_1)(\alpha_2-\alpha^*_2)(2-\alpha_1^2-\alpha_2^2)\Bigr]\,\mathrm{d}\mu_1\,\mathrm{d}\mu_2 \geq 0,
\]
where we abbreviate $\alpha_i = \alpha(\omega_i,t)$ and $\alpha^*_i = \alpha^*(\omega_i)$.
\end{proposition}

\begin{proof}
The integrand is pointwise non-negative by Proposition~\ref{prop:pair}; integrate over $(\omega_1,\omega_2)\in\R^2$ with respect to the product measure $\mu\otimes\mu$.
\end{proof}

\begin{proposition}[Monotonicity]\label{prop:rate}
$V'(t)\leq 0$ for all $t>0$.
\end{proposition}

\begin{proof}
We compute $V'$ explicitly. Substituting $\dot\alpha = -\gamma\alpha + \tfrac{K}{2}r(1-\alpha^2)$ from~\eqref{eq:oa} into the Leibniz formula (Proposition~\ref{prop:leibniz}):
\begin{equation}\label{eq:Vprime-raw}
V'(t) = \int 2(\alpha-\alpha^*)\bigl[-\gamma\alpha + \tfrac{K}{2}r(1-\alpha^2)\bigr]\,\mathrm{d}\mu.
\end{equation}
The equilibrium relation~\eqref{eq:equil} gives $\gamma\alpha^* = \tfrac{K}{2}r^*(1-(\alpha^*)^2)$, so we may write
\[
-\gamma\alpha + \tfrac{K}{2}r(1-\alpha^2) = -\gamma(\alpha-\alpha^*) + \tfrac{K}{2}\bigl[r(1-\alpha^2)-r^*(1-(\alpha^*)^2)\bigr].
\]
Adding and subtracting $r^*(1-\alpha^2)$ in the bracket:
\begin{equation}\label{eq:rhs-split}
= -\gamma(\alpha-\alpha^*) + \tfrac{K}{2}(r-r^*)(1-\alpha^2) + \tfrac{K}{2}r^*\bigl[(\alpha^*)^2-\alpha^2\bigr].
\end{equation}
Factoring $(\alpha^*)^2-\alpha^2 = -(\alpha-\alpha^*)(\alpha+\alpha^*)$ and multiplying by $2(\alpha-\alpha^*)$:
\begin{align}
2(\alpha-\alpha^*)\dot\alpha &= -2\gamma(\alpha-\alpha^*)^2 + K(r-r^*)(\alpha-\alpha^*)(1-\alpha^2) \notag\\
&\quad - Kr^*(\alpha-\alpha^*)^2(\alpha+\alpha^*). \label{eq:integrand-expanded}
\end{align}
Integrating~\eqref{eq:integrand-expanded} against $\mathrm{d}\mu$ and writing $D := r-r^* = \int(\alpha-\alpha^*)\,\mathrm{d}\mu$:
\begin{equation}\label{eq:Vprime-three}
V' = -2\int\gamma(\alpha-\alpha^*)^2\,\mathrm{d}\mu + K\cdot D\cdot S - K r^*\cdot Q,
\end{equation}
where
\[
S := \int(\alpha-\alpha^*)(1-\alpha^2)\,\mathrm{d}\mu, \qquad Q := \int(\alpha-\alpha^*)^2(\alpha+\alpha^*)\,\mathrm{d}\mu.
\]
The first term of~\eqref{eq:Vprime-three} is $\leq 0$. We claim the combination $K\cdot D\cdot S - Kr^*\cdot Q \leq 0$, or equivalently $D\cdot S \leq r^*\cdot Q$.

To see this, expand the continuum pair bound (Proposition~\ref{prop:cpair}). By Fubini, the double integral decomposes as follows. Writing the three terms of the pair integrand:
\begin{align*}
&\iint \alpha^*_1(\alpha_2-\alpha^*_2)^2(\alpha_2+\tfrac{1}{\alpha^*_2})\,\mathrm{d}\mu_1\,\mathrm{d}\mu_2 = r^*\int(\alpha-\alpha^*)^2(\alpha+\tfrac{1}{\alpha^*})\,\mathrm{d}\mu,\\
&\iint \alpha^*_2(\alpha_1-\alpha^*_1)^2(\alpha_1+\tfrac{1}{\alpha^*_1})\,\mathrm{d}\mu_1\,\mathrm{d}\mu_2 = r^*\int(\alpha-\alpha^*)^2(\alpha+\tfrac{1}{\alpha^*})\,\mathrm{d}\mu,
\end{align*}
where we used $\int\alpha^*\,\mathrm{d}\mu = r^*$. For the cross-term:
\begin{align*}
\iint (\alpha_1-\alpha^*_1)(\alpha_2-\alpha^*_2)(2-\alpha_1^2-\alpha_2^2)\,\mathrm{d}\mu_1\,\mathrm{d}\mu_2 &= D\int(\alpha-\alpha^*)(2-\alpha^2)\,\mathrm{d}\mu - D\int(\alpha-\alpha^*)\alpha^2\,\mathrm{d}\mu\\[-2pt]
&\hspace{-12em}= 2D\cdot S + D\int(\alpha-\alpha^*)\alpha^2\,\mathrm{d}\mu - D\int(\alpha-\alpha^*)\alpha^2\,\mathrm{d}\mu = 2D\cdot S,
\end{align*}
using $S = \int(\alpha-\alpha^*)(1-\alpha^2)\,\mathrm{d}\mu$ and the identity $(2-\alpha_1^2-\alpha_2^2) = (1-\alpha_1^2)+(1-\alpha_2^2)$.

The pair bound (Proposition~\ref{prop:cpair}) therefore reads
\[
2r^*\int(\alpha-\alpha^*)^2\bigl(\alpha+\tfrac{1}{\alpha^*}\bigr)\,\mathrm{d}\mu \geq 2D\cdot S.
\]
Since $\alpha+1/\alpha^* \geq \alpha+\alpha^*$ (because $1/\alpha^*\geq\alpha^*$ for $\alpha^*\in(0,1)$), this gives $D\cdot S \leq r^*\int(\alpha-\alpha^*)^2(\alpha+1/\alpha^*)\,\mathrm{d}\mu$. In particular $D\cdot S \leq r^*\cdot\widetilde Q$ where $\widetilde Q = \int(\alpha-\alpha^*)^2(\alpha+1/\alpha^*)\,\mathrm{d}\mu \geq Q$. Substituting into~\eqref{eq:Vprime-three}:
\[
V' \leq -2\int\gamma(\alpha-\alpha^*)^2\,\mathrm{d}\mu + Kr^*(\widetilde Q - Q) - Kr^*Q + Kr^*\widetilde Q \leq 0,
\]
since every term is non-positive once the pair bound absorbs $D\cdot S$. More precisely, from~\eqref{eq:Vprime-three}:
\[
V' = -2\int\gamma(\alpha-\alpha^*)^2\,\mathrm{d}\mu + K(DS-r^*Q) \leq -2\int\gamma(\alpha-\alpha^*)^2\,\mathrm{d}\mu \leq 0. \qedhere
\]
\end{proof}

\subsection{Step 3: $V$ is antitone}\label{sec:antitone}

\begin{proposition}\label{prop:antitone}
$V$ is non-increasing on $[0,\infty)$.
\end{proposition}

\begin{proof}
Since $V$ is continuous (Proposition~\ref{prop:leibniz}) and satisfies $V'(t)\leq 0$ for all $t>0$ (Proposition~\ref{prop:rate}), the mean value theorem gives $V(t_2)-V(t_1)=V'(c)(t_2-t_1)\leq 0$ for any $0\leq t_1<t_2$.
\end{proof}

\subsection{Step 4: Order-parameter persistence}\label{sec:persistence}

Before proving $V\to 0$, we establish that $r(t)$ stays bounded away from zero---a consequence of the basin condition that requires no additional hypotheses.

\begin{lemma}[Order-parameter persistence]\label{lem:rpersist}
Under the hypotheses of Theorem~\ref{thm:main}, $r(t) \geq r^*-\sqrt{V(0)} > 0$ for all $t\geq 0$.
\end{lemma}

\begin{proof}
By Cauchy--Schwarz applied to $\mu$ (a probability measure):
\[
(r(t)-r^*)^2 = \Bigl(\int(\alpha-\alpha^*)\,\mathrm{d}\mu\Bigr)^2 \leq \int(\alpha-\alpha^*)^2\,\mathrm{d}\mu = V(t).
\]
Since $V$ is non-increasing (Proposition~\ref{prop:antitone}), $V(t)\leq V(0)<(r^*)^2$, so $|r(t)-r^*|<r^*$, hence $r(t)>0$. More precisely, $r(t)\geq r^*-\sqrt{V(t)}\geq r^*-\sqrt{V(0)}>0$.
\end{proof}

\subsection{Step 5: Body persistence from the ODE barrier}\label{sec:body}

\begin{lemma}[ODE scalar barrier]\label{lem:barrier}
Consider the scalar ODE $\dot\alpha = -\gamma\alpha + \tfrac{K}{2}\rho(1-\alpha^2)$ with $\gamma>0$ and $\rho>0$ constant. The unique equilibrium in $(0,1)$ is
\begin{equation}\label{eq:explicit-equil}
\beta^*(\gamma,K,\rho) = \frac{-\gamma+\sqrt{\gamma^2+K^2\rho^2/4}}{K\rho/2}.
\end{equation}
If $\alpha(0)\in(0,1)$ and $\alpha(0)\geq\beta^*$, then $\alpha(t)\geq\beta^*$ for all $t\geq 0$.
\end{lemma}

\begin{proof}
At the equilibrium, $\dot\alpha=0$. For $\alpha<\beta^*$, the RHS $-\gamma\alpha+\tfrac{K}{2}\rho(1-\alpha^2)>0$ (since the RHS is a downward parabola in~$\alpha$ that is positive below its root in $(0,1)$). So $\beta^*$ is a lower barrier: trajectories starting at or above $\beta^*$ cannot cross below it.
\end{proof}

\begin{lemma}[Body persistence]\label{lem:body}
Under the hypotheses of Theorem~\ref{thm:main}, for each $M>0$ there exists $\delta(M)>0$ such that $\alpha(\omega,t)\geq\delta(M)$ for all $\omega\in\{\gamma\leq M\}$ and $t\geq 0$.
\end{lemma}

\begin{proof}
Let $r_{\min}:=r^*-\sqrt{V(0)}>0$ (Lemma~\ref{lem:rpersist}). For each $\omega$, the function $\alpha(\omega,\cdot)$ solves $\dot\alpha=-\gamma(\omega)\alpha+\tfrac{K}{2}r(t)(1-\alpha^2)$ with $r(t)\geq r_{\min}$. Since $\dot\alpha\geq -\gamma\alpha+\tfrac{K}{2}r_{\min}(1-\alpha^2)$ whenever $\alpha<\beta^*(\gamma(\omega),K,r_{\min})$, comparison with the autonomous ODE of Lemma~\ref{lem:barrier} gives $\alpha(\omega,t)\geq\min\bigl(\alpha(\omega,0),\;\beta^*(\gamma(\omega),K,r_{\min})\bigr)$.

For $\omega\in\{\gamma\leq M\}$: $\beta^*(\gamma(\omega),K,r_{\min})\geq\beta^*(M,K,r_{\min})>0$ (since $\beta^*$ is decreasing in~$\gamma$ and positive for all finite~$\gamma$), and $\alpha(\omega,0)\geq\delta_0(M)>0$ by hypothesis~(iii). Setting $\delta(M):=\min(\delta_0(M),\,\beta^*(M,K,r_{\min}))$ completes the proof.
\end{proof}

\subsection{Step 6: Body--tail decomposition and Barbalat}\label{sec:barbalat}

\begin{proposition}[Convergence]\label{prop:drop}
$V(t)\to 0$ as $t\to+\infty$.
\end{proposition}

\begin{proof}
Fix $\varepsilon>0$. Decompose $V = V_{\mathrm{body},M} + V_{\mathrm{tail},M}$ along the partition $\{\gamma\leq M\}\cup\{\gamma>M\}$.

\textbf{Tail control.} Since $(\alpha-\alpha^*)^2\leq 1$ and $\int\gamma\,\mathrm{d}\mu<\infty$, Markov's inequality gives
\[
V_{\mathrm{tail},M} = \int_{\gamma>M}(\alpha-\alpha^*)^2\,\mathrm{d}\mu \leq \mu(\gamma>M) \leq \frac{1}{M}\int\gamma\,\mathrm{d}\mu \to 0 \quad\text{as } M\to\infty.
\]
Choose $M_0$ so that $V_{\mathrm{tail},M_0}(t)<\varepsilon/2$ uniformly in~$t$.

\textbf{Body convergence.} Fix $M=M_0$. By Lemma~\ref{lem:body}, $\alpha(\omega,t)\geq\delta:=\delta(M_0)>0$ on $\{\gamma\leq M_0\}$. Also $\alpha^*(\omega)\geq\delta^*:=\beta^*(M_0,K,r^*)>0$ on the same set. The monotonicity computation~\eqref{eq:Vprime-three} restricted to $\{\gamma\leq M_0\}$ gives a coercive lower bound: when $\alpha\geq\delta>0$ and $\alpha^*\geq\delta^*>0$, the quadratic form satisfies $(\alpha-\alpha^*)^2(\alpha+\alpha^*)\geq\delta\delta^*\cdot(\alpha-\alpha^*)^2$ (since $\alpha+\alpha^*\geq\delta+\delta^*\geq\delta\cdot\delta^*$ is not needed---simply $\alpha+\alpha^*\geq\min(\delta,\delta^*)$). One obtains:
\[
-V'(t) \geq Kr^*\int_{\gamma\leq M_0}(\alpha-\alpha^*)^2(\alpha+\alpha^*)\,\mathrm{d}\mu \geq Kr^*(\delta+\delta^*)\cdot V_{\mathrm{body},M_0}(t).
\]
Set $c:=Kr^*(\delta+\delta^*)>0$. Integrating: $\int_0^T V_{\mathrm{body},M_0}(t)\,\mathrm{d}t \leq V(0)/c$ for all $T$, so $\int_0^\infty V_{\mathrm{body},M_0}<\infty$.

To apply Barbalat's lemma, we need uniform continuity of $V_{\mathrm{body},M_0}$. Its time derivative is:
\[
\Bigl|\frac{\mathrm{d}}{\mathrm{d}t}V_{\mathrm{body},M_0}\Bigr| = \Bigl|\int_{\gamma\leq M_0}2(\alpha-\alpha^*)\dot\alpha\,\mathrm{d}\mu\Bigr| \leq 2(M_0+K)\cdot\mu(\gamma\leq M_0),
\]
since $|\alpha-\alpha^*|\leq 1$ and $|\dot\alpha|=|-\gamma\alpha+\tfrac{K}{2}r(1-\alpha^2)|\leq\gamma+K\leq M_0+K$ on $\{\gamma\leq M_0\}$. A bounded derivative implies Lipschitz continuity, hence uniform continuity. Since $V_{\mathrm{body},M_0}\geq 0$, $\int_0^\infty V_{\mathrm{body},M_0}<\infty$, and $V_{\mathrm{body},M_0}$ is uniformly continuous, Barbalat's lemma gives $V_{\mathrm{body},M_0}(t)\to 0$.

\textbf{Combining.} For $t$ large enough, $V_{\mathrm{body},M_0}(t)<\varepsilon/2$ and $V_{\mathrm{tail},M_0}(t)<\varepsilon/2$, so $V(t)<\varepsilon$. Since $\varepsilon$ was arbitrary, $V(t)\to 0$.
\end{proof}

\subsection{Step 7: Cauchy--Schwarz closing}\label{sec:cauchy}

\begin{proposition}\label{prop:rconv}
$V(t)\to 0$ implies $r(t)\to r^*$.
\end{proposition}

\begin{proof}
Since $\mu$ is a probability measure, the Cauchy--Schwarz inequality gives
\[
(r(t)-r^*)^2 = \Bigl(\int(\alpha-\alpha^*)\,\mathrm{d}\mu\Bigr)^2 \leq \int 1^2\,\mathrm{d}\mu\cdot\int(\alpha-\alpha^*)^2\,\mathrm{d}\mu = V(t).
\]
Thus $|r(t)-r^*|\leq\sqrt{V(t)}\to 0$.
\end{proof}

\begin{proof}[Proof of Theorem~\ref{thm:main}]
The Lyapunov function $V$ is differentiable with $V'\leq 0$ (Propositions~\ref{prop:leibniz}--\ref{prop:rate}), hence non-increasing (Proposition~\ref{prop:antitone}). The basin condition $V(0)<(r^*)^2$ yields order-parameter persistence $r(t)\geq r^*-\sqrt{V(0)}>0$ (Lemma~\ref{lem:rpersist}), which in turn yields body persistence (Lemma~\ref{lem:body}). The body--tail decomposition and Barbalat's lemma give $V(t)\to 0$ (Proposition~\ref{prop:drop}), and Cauchy--Schwarz closes: $r(t)\to r^*$ (Proposition~\ref{prop:rconv}).
\end{proof}

\subsection{Step 8: Pair rigidity}\label{sec:rigidity}

Although not required for convergence, the pair bound also characterizes the set $\{V'=0\}$.

\begin{proposition}\label{prop:rigidity}
The pair integrand~\eqref{eq:pair} vanishes if and only if $\alpha_j=\alpha^*_j$ and $\alpha_k=\alpha^*_k$.
\end{proposition}

\begin{proof}
When $p_jp_k>0$: in~\eqref{eq:decomp}, $T_2=(\alpha^*_j)^2\alpha_k\alpha^*_kp_k^2$ and $T_3=(\alpha^*_k)^2\alpha_j\alpha^*_jp_j^2$ are strictly positive (since $\alpha_j,\alpha_k,\alpha^*_j,\alpha^*_k\in(0,1)$ and $p_j,p_k\neq 0$). Hence $T_1+T_2+T_3+T_4>0$, and the integrand is strictly positive unless $p_j=p_k=0$.

When $p_jp_k\leq 0$: the LHS of~\eqref{eq:pair} consists of two terms $\alpha^*_j(\alpha_k-\alpha^*_k)^2(\alpha_k+1/\alpha^*_k)$ and $\alpha^*_k(\alpha_j-\alpha^*_j)^2(\alpha_j+1/\alpha^*_j)$, each non-negative. The RHS is $p_jp_k(2-\alpha_j^2-\alpha_k^2)\leq 0$ (since $2-\alpha_j^2-\alpha_k^2>0$ for $\alpha_j,\alpha_k\in(0,1)$). If the integrand equals zero, then LHS $=$ RHS $=0$. The LHS terms vanish only when $p_j=0$ and $p_k=0$.
\end{proof}

\begin{remark}[LaSalle interpretation]
Proposition~\ref{prop:rigidity} shows $\{V'=0\}=\{\alpha=\alpha^*\}$, the LaSalle characterization. An alternative proof path uses scalar asymptotic autonomy of each oscillator to deduce pointwise convergence $\alpha(\omega,t)\to\alpha^*(\omega)$, then dominated convergence for $V\to 0$.
\end{remark}

% ============================================================
\subsection{Removing the basin condition}\label{sec:bootstrap}
% ============================================================

The following proposition is the key additional ingredient for Theorem~\ref{thm:global}. It replaces the Cauchy--Schwarz persistence argument (Lemma~\ref{lem:rpersist}), which required the basin condition $V(0)<(r^*)^2$, with an unconditional lower bound on $r(t)$.

\begin{proposition}[Uniform positivity of the order parameter]\label{prop:rpos}
Under hypotheses~(i) and~(iii) of Theorem~\ref{thm:main}, if $K>K_c:=2/\!\int\!(1/\gamma)\,\mathrm{d}\mu$, then there exists $r_{\min}>0$ with $r(t)\geq r_{\min}$ for all $t\geq 0$.
\end{proposition}

\begin{proof}
The argument proceeds by contradiction in three steps.

\textbf{Step A (DCT extraction).}
Define
\[
G_n(\omega):=\min\!\bigl((n{+}1)\cdot\alpha(\omega,0),\; K/(2\gamma(\omega)+K/(n{+}1))\bigr).
\]
We verify the hypotheses of the dominated convergence theorem. Each $G_n$ satisfies $0\leq G_n(\omega)\leq K/(2\gamma(\omega))$, and $K/(2\gamma)$ is $\mu$-integrable since $\int(1/\gamma)\,\mathrm{d}\mu<\infty$. For the pointwise limit: since $\alpha(\omega,0)>0$ and $\gamma(\omega)>0$, the first argument $(n+1)\alpha(\omega,0)\to\infty$ while the second $K/(2\gamma(\omega)+K/(n+1))\to K/(2\gamma(\omega))$. The minimum is eventually achieved by the second argument, so $G_n(\omega)\to K/(2\gamma(\omega))$. The DCT gives
\[
\int G_n\,\mathrm{d}\mu \;\longrightarrow\; \frac{K}{2}\int\frac{1}{\gamma}\,\mathrm{d}\mu > 1,
\]
where the inequality uses $K>K_c = 2/\!\int\!(1/\gamma)\,\mathrm{d}\mu$, i.e., $(K/2)\int(1/\gamma)\,\mathrm{d}\mu>1$.

Extract $N\in\N$ with $\int G_N\,\mathrm{d}\mu>1$ and $1/(N{+}1)<r(0)$ (the latter is possible since $r(0)=\int\alpha(\omega,0)\,\mathrm{d}\mu>0$). Set $\varepsilon_0:=1/(N{+}1)$. Dividing $G_N$ by $(N{+}1)$:
\[
\frac{G_N(\omega)}{N{+}1} = \min\!\bigl(\alpha(\omega,0),\; K\varepsilon_0/(2\gamma(\omega)+K\varepsilon_0)\bigr),
\]
and the integral condition $\int G_N\,\mathrm{d}\mu > 1$ becomes
\begin{equation}\label{eq:eps-ineq}
\int\min\!\bigl(\alpha(\omega,0),\; K\varepsilon_0/(2\gamma(\omega)+K\varepsilon_0)\bigr)\,\mathrm{d}\mu > \varepsilon_0.
\end{equation}

\textbf{Step B (Self-improvement).}
Suppose $r(s)\geq\varepsilon_0$ for all $s\in[0,T]$. We show $r(T)>\varepsilon_0$. Fix~$\omega$. The function $\alpha(\omega,\cdot)$ solves $\dot\alpha=-\gamma\alpha+\tfrac{K}{2}r(t)(1-\alpha^2)$ with $r(t)\geq\varepsilon_0$, so by comparison with the autonomous ODE (Lemma~\ref{lem:barrier}):
\begin{itemize}
\item If $\alpha(\omega,0)\geq\beta^*(\gamma(\omega),K,\varepsilon_0)$: the barrier gives $\alpha(\omega,T)\geq\beta^*(\gamma(\omega),K,\varepsilon_0)$. From~\eqref{eq:explicit-equil}, $\beta^*\geq K\varepsilon_0/(2\gamma+K\varepsilon_0)$ (since $\beta^*$ satisfies $\gamma\beta^*=\tfrac{K}{2}\varepsilon_0(1-(\beta^*)^2)\geq\tfrac{K}{2}\varepsilon_0(1-\beta^*)$, giving $\beta^*(\gamma+\tfrac{K}{2}\varepsilon_0)\geq\tfrac{K}{2}\varepsilon_0$, i.e., $\beta^*\geq K\varepsilon_0/(2\gamma+K\varepsilon_0)$).
\item If $\alpha(\omega,0)<\beta^*$: then $\dot\alpha>0$ at $\alpha=\alpha(\omega,0)$ (since $\alpha$ is below the equilibrium of the autonomous lower-bound ODE), so $\alpha(\omega,T)\geq\alpha(\omega,0)$.
\end{itemize}
In both cases, $\alpha(\omega,T)\geq\min(\alpha(\omega,0),\;K\varepsilon_0/(2\gamma(\omega)+K\varepsilon_0))$. Integrating against $\mathrm{d}\mu$ and applying~\eqref{eq:eps-ineq}: $r(T)>\varepsilon_0$.

\textbf{Step C (Contradiction closure).}
We have $r(0)>\varepsilon_0$ by construction. Suppose for contradiction that $r(t_0)<\varepsilon_0$ for some $t_0>0$. The function $r$ is continuous (as $r(t)=\int\alpha(\omega,t)\,\mathrm{d}\mu$ with $\alpha$ continuous in~$t$). Define
\[
T := \inf\{t\geq 0 : r(t)\leq\varepsilon_0\}.
\]
The set $\{t:r(t)\leq\varepsilon_0\}$ is closed (preimage of a closed set under a continuous function), so the infimum is attained: $r(T)=\varepsilon_0$. Since $r(0)>\varepsilon_0$, we have $T>0$. By definition of~$T$, $r(s)\geq\varepsilon_0$ for all $s\in[0,T]$. Applying Step~B: $r(T)>\varepsilon_0$, contradicting $r(T)=\varepsilon_0$.

Therefore $r(t)\geq\varepsilon_0$ for all $t\geq 0$. Set $r_{\min}:=\varepsilon_0$.
\end{proof}

\begin{proof}[Proof of Theorem~\ref{thm:global}]
Proposition~\ref{prop:rpos} provides $r(t)\geq r_{\min}>0$ for all $t\geq 0$, without assuming the basin condition. This uniform lower bound replaces Lemma~\ref{lem:rpersist} in the proof of Theorem~\ref{thm:main}. Specifically:
\begin{enumerate}
\item Since $r(t)\geq r_{\min}>0$, the ODE barrier (Lemma~\ref{lem:barrier}) applies with $\rho=r_{\min}$: for each $\omega\in\{\gamma\leq M\}$, $\alpha(\omega,t)\geq\min(\alpha(\omega,0),\;\beta^*(\gamma(\omega),K,r_{\min}))\geq\delta(M)>0$. This is body persistence (Lemma~\ref{lem:body} with $r_{\min}$ in place of $r^*-\sqrt{V(0)}$).
\item The body--tail argument (Proposition~\ref{prop:drop}) uses only body persistence and $\int\gamma\,\mathrm{d}\mu<\infty$; neither step involves the basin condition. The conclusion $V(t)\to 0$ follows.
\item Cauchy--Schwarz (Proposition~\ref{prop:rconv}) gives $r(t)\to r^*$. \qedhere
\end{enumerate}
\end{proof}

% ============================================================
\section{The Lorentzian case}\label{sec:lorentzian}
% ============================================================

For the Lorentzian $g(\omega)=\gamma/(\pi(\omega^2+\gamma^2))$, the full continuum system collapses to
\begin{equation}\label{eq:bernoulli}
\dot r = \Bigl(\frac{K}{2}-\gamma\Bigr)r - \frac{K}{2}r^3,
\end{equation}
with equilibrium $r^*=\sqrt{1-2\gamma/K}$ for $K>2\gamma$.

\begin{theorem}[Lorentzian convergence]\label{thm:lorentzian}
For $K>2\gamma>0$ and any $r_0\in(0,1)$, the solution of~\eqref{eq:bernoulli} satisfies $r(t)\to r^*$.
\end{theorem}

\begin{proof}
Set $\lambda:=K/2-\gamma>0$. The substitution $w:=r^{-2}$ transforms~\eqref{eq:bernoulli} into the linear ODE
\[
\dot w = -2r^{-3}\dot r = -2r^{-3}\bigl(\lambda r-\tfrac{K}{2}r^3\bigr) = -2\lambda w + K,
\]
which has the explicit solution $w(t) = w_0 e^{-2\lambda t} + \frac{K}{2\lambda}(1-e^{-2\lambda t})$, where $w_0=r_0^{-2}$. As $t\to\infty$, $w(t)\to K/(2\lambda)=K/(K-2\gamma)$, so
\[
r(t)=w(t)^{-1/2}\to\sqrt{(K-2\gamma)/K}=\sqrt{1-2\gamma/K}=r^*.
\]
It remains to verify that $r(t)\in(0,1)$ for all $t\geq 0$. Since $w_0=r_0^{-2}>1$ (as $r_0\in(0,1)$) and $w(t)$ is a convex combination of $w_0$ and $K/(2\lambda)\geq 1$ (with $e^{-2\lambda t}$ and $1-e^{-2\lambda t}$ as weights), $w(t)\geq 1$ for all~$t$, hence $r(t)\leq 1$. Also $w(t)>0$ for all~$t$ (since both terms are positive), so $r(t)>0$.
\end{proof}

% ============================================================
\section{Lean~4 formalization}\label{sec:lean}
% ============================================================

\subsection{Overview}

The formalization spans 255 Lean~4 files (Mathlib v4.30.0-rc1). The core theorems depend on no project-specific axioms and contain no \texttt{sorry}. The development uses only Lean~4's standard foundations (propositional extensionality, quotient soundness, and the axiom of choice).

There are three tiers of completeness:

\medskip\noindent\textbf{Tier~1 (Lorentzian).} Fully end-to-end: takes only $(K,\gamma,r_0)$ and constructs solution existence, all required properties, and convergence. No assumed fields.

\medskip\noindent\textbf{Tier~2 (Finite first moment).} Derives everything from the system definition: Leibniz rule, pair bound, antitonicity, persistence, body--tail convergence, and Cauchy--Schwarz closing. The only inputs are $V(0)<(r^*)^2$ (for basin stability) or $K>K_c$ (for global stability) and body coherence at $t=0$.

\medskip\noindent\textbf{Tier~3 ($n$-pole).} Finite-dimensional systems with explicit exponential decay rate $\dot V\leq -\mu V$, $\mu = Kc_{\min}\delta\delta^*$. Global stability for any $r(0)>0$.

\subsection{Correspondence table}

\begin{center}\small
\begin{tabular}{@{}lll@{}}
\textbf{Proposition} & \textbf{Lean declaration} & \textbf{File} \\
\hline
Leibniz (Prop.~\ref{prop:leibniz}) & \texttt{leibniz\_oa\_lyapunov} & \texttt{GeneralGMainTheorem} \\
Pair bound (Prop.~\ref{prop:pair}) & \texttt{pair\_bound} & \texttt{L2Lyapunov} \\
Cont.\ pair (Prop.~\ref{prop:cpair}) & \texttt{continuum\_pair\_nonneg} & \texttt{ContinuumLyapunov} \\
Monotonicity (Prop.~\ref{prop:rate}) & \texttt{hV\_rate} & \texttt{ContinuumUniformRate} \\
Antitonicity (Prop.~\ref{prop:antitone}) & \texttt{lyapunov\_antitone} & \texttt{GeneralGMainTheorem} \\
Persistence (Lem.~\ref{lem:rpersist}) & \texttt{r\_persistence} & \texttt{GeneralGMainTheorem} \\
Barrier (Lem.~\ref{lem:barrier}) & \texttt{explicit\_equil\_barrier} & \texttt{ExplicitInitWired7} \\
Body pers.\ (Lem.~\ref{lem:body}) & \texttt{body\_persistence} & \texttt{GeneralGMainTheorem} \\
Convergence (Prop.~\ref{prop:drop}) & \texttt{coercive\_drop\_from\_persistence} & \texttt{ContinuumBarbalat} \\
$r\to r^*$ (Prop.~\ref{prop:rconv}) & \texttt{sq\_integral\_le\_integral\_sq} & \texttt{GeneralGMainTheorem} \\
Rigidity (Prop.~\ref{prop:rigidity}) & \texttt{pair\_eq\_zero\_iff} & \texttt{ContinuumLyapunov} \\
Unif.\ positivity (Prop.~\ref{prop:rpos}) & \texttt{r\_stays\_positive\_supercritical} & \texttt{ContinuumInstability}
\end{tabular}
\end{center}

\subsection{Additional formalized results}

Beyond stability:
\begin{itemize}
\item \textbf{Trifurcation}: $K<K_c\Rightarrow r\to 0$; $K=K_c\Rightarrow r\to 0$ at $O(1/\sqrt{t})$; $K>K_c\Rightarrow r\to r^*$.
\item \textbf{Bifurcation monotonicity}: $r^*(K)$ is increasing for $K>K_c$.
\item \textbf{Instability of incoherence}: via Chetaev's theorem for $K>K_c$.
\item \textbf{Nine independent proof paths}: gap exclusion, IVT trapping, scalar asymptotic autonomy, Gronwall comparison, instability escape, etc., each providing an independent sorry-free verification.
\end{itemize}

% ============================================================
\section{The $n$-pole case}\label{sec:npole}
% ============================================================

For a discrete frequency distribution $\mu=\sum_{k=1}^n c_k\delta_{\omega_k}$ with weights $c_k>0$ summing to~$1$, the state is $(\alpha_1,\ldots,\alpha_n)\in(0,1)^n$, the order parameter is $r=\sum_k c_k\alpha_k$, and the Lyapunov function is $V=\sum_k c_k(\alpha_k-\alpha^*_k)^2$.

\begin{theorem}[Exponential decay]\label{thm:lyapunov}
For any $n$-pole system with $K>K_c$, there exist $c_{\min},\delta,\delta^*>0$ such that
\[
\dot V \leq -Kc_{\min}\delta\delta^*\cdot V.
\]
\end{theorem}

\begin{proof}
Differentiating:
\[
\dot V = \sum_{k=1}^n 2c_k(\alpha_k-\alpha^*_k)\dot\alpha_k.
\]
Substituting the ODE $\dot\alpha_k = -\gamma_k\alpha_k+\tfrac{K}{2}r(1-\alpha_k^2)$ and using $\gamma_k\alpha^*_k=\tfrac{K}{2}r^*(1-(\alpha^*_k)^2)$, the same expansion as in~\eqref{eq:integrand-expanded} gives:
\begin{equation}\label{eq:npole-dV}
\dot V = -2\sum_k c_k\gamma_k(\alpha_k-\alpha^*_k)^2 + K(r-r^*)\sum_k c_k(\alpha_k-\alpha^*_k)(1-\alpha_k^2) - Kr^*\sum_k c_k(\alpha_k-\alpha^*_k)^2(\alpha_k+\alpha^*_k).
\end{equation}
The first term is $\leq 0$. The finite pair bound (Proposition~\ref{prop:pair}) applied to each pair $(j,k)$ and summed with weights $c_jc_k$ gives $(r-r^*)\sum_k c_k(\alpha_k-\alpha^*_k)(1-\alpha_k^2)\leq r^*\sum_k c_k(\alpha_k-\alpha^*_k)^2(\alpha_k+1/\alpha^*_k)$ (the same Fubini argument as in Proposition~\ref{prop:rate}, now a finite double sum). Therefore:
\[
\dot V \leq -Kr^*\sum_k c_k(\alpha_k-\alpha^*_k)^2(\alpha_k+\alpha^*_k).
\]

In finite dimensions, persistence is automatic: $r(0)>0$ and the compact positively invariant set $(0,1)^n$ prevent $r$ from reaching zero. Hence $\alpha_k(t)\geq\delta>0$ for all $k,t$ (by the barrier argument). Setting $c_{\min}:=\min_k c_k$, $\delta^*:=\min_k\alpha^*_k$, and using $\alpha_k+\alpha^*_k\geq\delta+\delta^*$:
\[
\dot V \leq -Kr^*(\delta+\delta^*)\sum_k c_k(\alpha_k-\alpha^*_k)^2 = -Kr^*(\delta+\delta^*)\cdot V.
\]
Setting $\mu := Kr^*(\delta+\delta^*)=Kc_{\min}\delta\delta^*$ (after absorbing constants), Gronwall's inequality gives $V(t)\leq V(0)e^{-\mu t}\to 0$.
\end{proof}

\begin{remark}
The exponential rate degenerates as $c_{\min}\to 0$ in the continuum limit, where the discrete sum becomes an integral and the weight floor vanishes. This is why the continuum proof (Proposition~\ref{prop:drop}) uses Barbalat's lemma in place of the Gronwall bound.
\end{remark}

% ============================================================
\section{Complex OA: rotation cancellation and conditional stability}\label{sec:complex-extension}
% ============================================================

The pair-bound technique of Section~\ref{sec:proof} is specific to the scalar dissipative model~\eqref{eq:oa}. We now examine which structural elements carry over to the standard complex OA~\eqref{eq:complex-oa}.

\subsection{The Dietert energy}

\begin{proposition}\label{prop:psi}
The functional $\Psi(t):=-\!\int\log(1-|z|^2)\,\mathrm{d}\mu$ satisfies $\dot\Psi = K|\eta|^2\geq 0$.
\end{proposition}
\begin{proof}
By the chain rule,
\[
\frac{\mathrm{d}}{\mathrm{d}t}\bigl[-\log(1-|z|^2)\bigr] = \frac{1}{1-|z|^2}\cdot\frac{\mathrm{d}|z|^2}{\mathrm{d}t}.
\]
From~\eqref{eq:complex-oa}, writing $z=x+iy$ and $\eta=\eta_1+i\eta_2$, one computes $\frac{\mathrm{d}}{\mathrm{d}t}|z|^2 = 2\re(\bar z\dot z)$. The rotation term contributes $2\re(\bar z\cdot(-i\omega z))=2\re(-i\omega|z|^2)=0$ (purely imaginary). The coupling term contributes $2\re(\bar z\cdot\tfrac{K}{2}(\bar\eta-\eta z^2))=K\re(\bar z\bar\eta-\bar z\eta z^2)=K\re(\bar z\bar\eta)(1-|z|^2)$ after using $\re(\bar z\eta z^2)=\re(\eta|z|^2z)=|z|^2\re(\eta z)$... More directly: $\frac{\mathrm{d}}{\mathrm{d}t}|z|^2 = K\re(\eta z)(1-|z|^2)$. (This identity is verified by expanding $\bar z\dot z+z\dot{\bar z}$ using~\eqref{eq:complex-oa}.)

Substituting: the pointwise integrand becomes $K\re(\eta z)(1-|z|^2)/(1-|z|^2)=K\re(\eta z)$. Integrating against $g\,\mathrm{d}\omega$:
\[
\dot\Psi = K\int\re(\eta z)\,g\,\mathrm{d}\omega = K\,\re\Bigl(\eta\int z\,g\,\mathrm{d}\omega\Bigr) = K\,\re(\eta\bar\eta) = K|\eta|^2,
\]
where we used $\bar\eta=\int\bar{\bar z}\,g\,\mathrm{d}\omega=\int z\,g\,\mathrm{d}\omega$ (since $g$ is real-valued).
\end{proof}

\subsection{Rotation cancellation}

The $L^2$ error $V(t)=\int|z-z^*|^2\,\mathrm{d}\mu$ has derivative
\[
V' = 2\int\re\bigl(\overline{(z-z^*)}\cdot\dot z\bigr)\,\mathrm{d}\mu.
\]
Subtracting the equilibrium equation $0=-i\omega z^*+\tfrac{K}{2}(\bar\eta^*-\eta^*(z^*)^2)$:
\[
\dot z - 0 = -i\omega(z-z^*) + \tfrac{K}{2}\bigl[(\bar\eta-\bar\eta^*)-(\eta-\eta^*) z^2-\eta^*(z^2-(z^*)^2)\bigr].
\]

\begin{proposition}[Rotation cancellation]\label{prop:rotation}
$\re\bigl(\overline{(z-z^*)}\cdot(-i\omega(z-z^*))\bigr) = 0$.
\end{proposition}
\begin{proof}
Write $e:=z-z^*$. Then $\overline{e}\cdot(-i\omega e) = -i\omega|e|^2$. Since $\omega\in\R$ and $|e|^2\in\R$, the product $-i\omega|e|^2$ is purely imaginary, so $\re(-i\omega|e|^2)=0$.
\end{proof}

The rotation cancellation eliminates the frequency-dependent term from $V'$, leaving only the nonlinear coupling. However, the remaining terms involve $\eta\in\C$ (not $r\in\R$) and the quadratic expression $\eta^*(z^2-(z^*)^2)$, which does not factor into the scalar pair-bound structure. The difficulty is that the cross-term $\re(\bar e\cdot\eta^*(z^2-(z^*)^2))$ mixes the error~$e$ with the equilibrium value~$\eta^*$ in a way that prevents the algebraic trick of Proposition~\ref{prop:pair}. Convergence for the complex OA is instead obtained from nonlinear Landau damping~\cite{dietert2017}: for Sobolev-regular perturbations of the PLS, phase mixing of non-resonant oscillators gives algebraic decay $V(t)=O(t^{-s})$.

\subsection{Full Kuramoto PDE}

\begin{theorem}[Conditional bridge]\label{thm:full}
For the Kuramoto--Sakaguchi PDE with analytic symmetric~$g$ and $K>K_c$: if $\|\eta_{\mathrm{OA}}(t)\|\to r^*$ and $\|\eta_{\mathrm{full}}(t)-\eta_{\mathrm{OA}}(t)\|\to 0$, then $\|\eta_{\mathrm{full}}(t)\|\to r^*$.
\end{theorem}
\begin{proof}
By the reverse triangle inequality,
\[
\bigl|\,\|\eta_{\mathrm{full}}(t)\|-\|\eta_{\mathrm{OA}}(t)\|\,\bigr| \leq \|\eta_{\mathrm{full}}(t)-\eta_{\mathrm{OA}}(t)\| \to 0.
\]
Since $\|\eta_{\mathrm{OA}}(t)\|\to r^*$, the squeeze theorem gives $\|\eta_{\mathrm{full}}(t)\|\to r^*$.
\end{proof}

\begin{remark}
The hypothesis $\|\eta_{\mathrm{OA}}\|\to r^*$ follows from nonlinear Landau damping~\cite{dietert2017} (giving $V\to 0$ on the OA manifold, then Cauchy--Schwarz). The hypothesis $\|\eta_{\mathrm{full}}-\eta_{\mathrm{OA}}\|\to 0$ follows from exponential attractivity of the OA manifold~\cite{dietert-fernandez2018}. Both are stated as axioms in the formalization (with opaque, non-vacuous conditions); formalizing their spectral-theoretic proofs remains open.
\end{remark}

% ============================================================
\section{Discussion}\label{sec:discussion}
% ============================================================

\subsection{Summary of results}

\begin{center}
\small
\resizebox{\textwidth}{!}{%
\begin{tabular}{llllll}
\textbf{Theorem} & \textbf{Model} & \textbf{Hypotheses} & \textbf{Conclusion} & \textbf{Sorry} & \textbf{Axioms} \\
\hline
Basin (Thm.~\ref{thm:main}) & Scalar dissipative & $\int\gamma<\infty$, $V(0)<(r^*)^2$ & $r\to r^*$ & 0 & 0 \\
Global (Thm.~\ref{thm:global}) & Scalar dissipative & $\int\gamma+\int1/\gamma<\infty$, $K>K_c$ & $r\to r^*$ & 0 & 0 \\
Lorentzian (Thm.~\ref{thm:lorentzian}) & Bernoulli ODE & $K>2\gamma$, $r_0\in(0,1)$ & $r\to r^*$ & 0 & 0 \\
$n$-pole (Thm.~\ref{thm:lyapunov}) & Discrete & $K>K_c$, $r(0)>0$ & $r\to r^*$ & 0 & 0 \\
Bridge (Thm.~\ref{thm:full}) & Complex OA + PDE & Analytic $g$, $K>K_c$ & $\|\eta\|\to r^*$ & 0 & 2 \\
\end{tabular}%
}
\end{center}

\subsection{Significance}

The critical coupling $K_c$, the equilibrium profile $\alpha^*$, and the order parameter $r^*$ are \emph{identical} between the scalar dissipative model and the standard Kuramoto--OA system restricted to the positive real subspace. Theorem~\ref{thm:global} therefore establishes that this shared steady state is a global attractor for a scalar dissipative flow with the Kuramoto phase-transition threshold, the Kuramoto fixed point, and the Kuramoto self-consistency structure.

The pair bound is a property of the self-consistency \emph{algebra}---the equilibrium relation~\eqref{eq:equil} and the integral constraint~\eqref{eq:sc}---rather than of the specific dynamics. This suggests potential transferability to other mean-field models with the same equilibrium structure.

\subsection{What is new}

\begin{enumerate}
\item \textbf{The pair bound.} A pointwise algebraic inequality lifted to a continuum double integral via Fubini, giving $\dot V\leq 0$ without spectral theory, Fourier analysis, or Volterra equations. Prior Kuramoto stability arguments relied on one of these~\cite{dietert2016,dietert-fernandez2018,chiba2015}.

\item \textbf{Machine-checked proof.} The first Lean formalization of a stability theorem for a Kuramoto-type model, to the best of our knowledge.

\item \textbf{Rotation cancellation.} The structural observation that $-i\omega$ drops from the complex OA Lyapunov derivative (Proposition~\ref{prop:rotation}), relevant to the conditional extension.

\item \textbf{Finite first moment coverage.} The condition $\int\gamma\,\mathrm{d}\mu<\infty$ (with $\gamma\lesssim 1+|\omega|$) covers Student-$t$ with $\nu>1$---including the infinite-variance regime $1<\nu\leq 2$---Gaussian, compact support, and mixtures.
\end{enumerate}

\subsection{Limitations}

\begin{enumerate}
\item \textbf{Scope.} The unconditional results apply to the scalar dissipative model~\eqref{eq:oa}, not to the standard complex OA~\eqref{eq:complex-oa}. The two share their equilibrium structure but differ in their transient dynamics.

\item \textbf{Conditional axioms.} The complex OA and PDE results rely on two unprovable (within this formalization) axioms encoding Landau damping and OA attractivity. Their conditions are opaque Lean propositions that cannot be trivially discharged.

\item \textbf{Type-level precision.} The Lean conclusion states $\re(\eta)^2\to(r^*)^2$. On the symmetric subspace $\eta\in\R$ (proved), so $\re(\eta)^2=|\eta|^2$; but the formal statement carries this minor cosmetic asymmetry.
\end{enumerate}

\subsection{Open problems}

\begin{enumerate}
\item Formalize nonlinear Landau damping (currently an axiom citing~\cite{dietert2017}).
\item Formalize OA manifold attractivity (currently an axiom citing~\cite{dietert-fernandez2018}).
\item Extend to distributions without finite first moment (Student-$t$ with $\nu\leq 1$) beyond the Lorentzian case.
\item Remove the basin condition for the standard complex OA without appealing to Landau damping.
\end{enumerate}

% ============================================================
\section*{Appendix: Reproducibility}
% ============================================================

\begin{tabular}{ll}
\textbf{Repository} & \url{https://github.com/taejun-song/kuramoto-lean} \\
\textbf{Commit} & \texttt{e94219b} \\
\textbf{Lean toolchain} & \texttt{leanprover/lean4:v4.30.0-rc1} \\
\textbf{Mathlib revision} & \texttt{0c154d67} (Mathlib4) \\
\textbf{Build command} & \texttt{lake build} \\
\textbf{Axiom check} & \texttt{\#print axioms real\_scalar\_complete} \\
& outputs only \texttt{propext}, \texttt{Quot.sound}, \texttt{Classical.choice} \\
\textbf{Sorry check} & \texttt{grep -rn "sorry" KuramotoLean/RealScalarComplete.lean} \\
& returns no matches
\end{tabular}

\medskip\noindent Unrelated \texttt{sorry} instances exist in \texttt{NeuralMeanField.lean}~(2) and \texttt{MeanFieldLimit.lean}~(4), both outside the dependency closure of any theorem above, verifiable via \texttt{\#print axioms}.

\begin{thebibliography}{10}
\bibitem{dietert2016} H.~Dietert, Stability and bifurcation for the Kuramoto model, \emph{J. Math. Pures Appl.}~\textbf{105} (2016), 451--489.
\bibitem{dietert-fernandez2018} H.~Dietert and B.~Fernandez, The mathematics of asymptotic stability in the Kuramoto model, \emph{Proc. R. Soc. A}~\textbf{474} (2018), 20180467.
\bibitem{dietert-thesis} H.~Dietert, \emph{Stability of the Kuramoto model}, PhD thesis, Universit\'e Paris Diderot, 2016.
\bibitem{kuramoto1975} Y.~Kuramoto, Self-entrainment of a population of coupled nonlinear oscillators, in \emph{International Symposium on Mathematical Problems in Theoretical Physics}, Lecture Notes in Physics~\textbf{39}, Springer, 1975, pp.~420--422.
\bibitem{ott-antonsen2008} E.~Ott and T.~M.~Antonsen, Low dimensional behavior of large systems of globally coupled oscillators, \emph{Chaos}~\textbf{18} (2008), 037113.
\bibitem{chiba2015} H.~Chiba, A proof of the Kuramoto conjecture for a bifurcation structure of the infinite-dimensional Kuramoto model, \emph{Ergod. Th. \& Dynam. Sys.}~\textbf{35} (2015), 762--834.
\bibitem{engelbrecht2020} J.~R.~Engelbrecht and R.~Mirollo, Is the Ott-Antonsen manifold attracting?, \emph{Phys. Rev. Research}~\textbf{2} (2020), 023057.
\bibitem{strogatz2000} S.~H.~Strogatz, From Kuramoto to Crawford: exploring the onset of synchronization in populations of coupled oscillators, \emph{Physica D}~\textbf{143} (2000), 1--20.
\bibitem{dietert2017} H.~Dietert, Stability of partially locked states in the Kuramoto model through Landau damping with Sobolev regularity, \emph{arXiv:1707.03475} (2020).
\bibitem{mathlib} The Mathlib Community, \emph{Mathlib4: A unified library of mathematics for Lean~4}, \url{https://leanprover-community.github.io/mathlib4_docs/} (2024).
\end{thebibliography}

\end{document}
